\subsection{Computational Expressiveness of Splicing}
\label{annex:p09-recombination-computation}

\begin{annexitem}{annex:p09.recombination-regular-simulation}{thm:p09-recombination-regular-simulation}{Regular rules, complete simulation, and terminal extraction}
\begin{proof}
Each rule group is a finite union of fixed words with one occurrence of
\(U^*\) or \(T^*\) substituted at the displayed position. A finite
automaton reads the fixed prefix, loops over the indicated finite alphabet,
and reads the fixed suffix, nondeterministically selecting the split if
necessary. Finite union proves regularity of \(R\). The seed is finite.

Call \(XbY\), with \(b\in U^*\) containing exactly one \(B\), an ordinary
active word. Initially it is \(XBSY\), decoding to \(S\).
For \(b=pBq\) define its decoding to be \(qp\).
Groups 2--5 give, by direct substitution into the four rule handles,
\[
\begin{aligned}
 (Xw\alpha Y,ZY_\alpha)&\vdash XwY_\alpha,\\
 (X'\alpha Z,XwY_\alpha)&\vdash X'\alpha wY_\alpha,\\
 (X'\alpha wY_\alpha,ZY)&\vdash X'\alpha wY,\\
 (XZ,X'\alpha wY)&\vdash X\alpha wY.
\end{aligned}
\]
The first step stores the removed letter in \(Y_\alpha\); the second
restores exactly that letter. No two different stored letters can match.
Repeating the four steps gives all circular rotations of \(b\).
There is still exactly one \(B\) at every completed phase, and the
decoding is unchanged, including a rotation in which \(\alpha=B\).

If an occurrence \(u\) is chosen in the decoded sentential form, rotate
the circle to place the gap just after that occurrence at its right end.
Since \(u\) contains no \(B\), it does not cross the marked beginning.
The resulting word is \(XwuY\) for some \(w\in U^*\).
Group 1 with helper \(ZvY\) gives \(XwvY\). This changes the decoded
form by exactly the chosen grammar production \(u\to v\).
Rotate again to put \(B\) first. Induction on grammar derivation length
therefore gives \(XBqY\) for every sentential form \(q\) of \(G\).

For the converse, classify every generated word. Besides the unchanged
helpers, there are ordinary active words, the three unfinished rotation
phases in the display, extraction words \(qY\) with \(q\in T^*\), their
terminal descendants \(q\), and inert words formed by prematurely
removing the final \(Y\). One additional inert helper class is \(Zv\):
group 7 on \(ZvY\) retains \(Z\). This is why no general assertion
that outputs contain no \(Z\) is needed.

Here is the complete check on the classes. Group 1 requires an initial
\(X\), a suffix \(uY\), and a helper \(ZvY\). Its first operand must
be an ordinary active word; \(X'\)-phases and \(Y_\alpha\)-phases
cannot supply its sites. Group 2 likewise requires an ordinary word
ending in \(\alpha Y\), and its second site \(ZY_\alpha\) selects
the corresponding helper. Group 3 requires the helper \(X'\alpha Z\)
and a phase beginning in \(X\) and ending in \(Y_\alpha\).
Group 4 requires an \(X'\)-phase ending in \(Y_\alpha\) and helper
\(ZY\). Group 5 requires \(XZ\) and an \(X'\)-phase ending in
\(wY\). No other listed word supplies those complete recognition sites.
Thus these groups produce exactly the next displayed phase or an ordinary
word obtained by a grammar rewrite or rotation.

Group 6 requires the helper \(ZY\) and an ordinary word \(XBqY\)
whose whole suffix \(q\) lies in \(T^*\); it produces \(qY\).
Group 7 can remove a final \(Y\) from an ordinary word, an \(X'\)-phase,
an extraction word, or a \(ZvY\) helper. In the first two cases the
remaining \(X\) or \(X'\) cannot be removed: groups 1--6 require a
trailing \(Y\) or \(Y_\alpha\), or a trailing \(Z\) in the particular
helper form. These prematurely closed words have none. Words \(Zv\)
are also inert because \(v\in(N\cup T)^*\), so they have no \(Y\)
or tagged symbol and do not supply \(XZ\) or \(X'\alpha Z\).
The only terminal output of group 7 is therefore the \(q\) descended
from \(qY\).

This class check is an induction on set-splicing stages; it also proves
that the decoding of every ordinary active word is a grammar sentential
form. For \(q\in T^*\) that decoding is a generated terminal word.
No helper lies in \(T^*\). Hence every terminal output is in \(L(G)\).
Conversely, the grammar simulation produces \(XBqY\) for every
\(q\in L(G)\), then groups 6 and 7 produce \(q\).
If \(q=\lambda\), these steps produce \(Y\) and then \(\lambda\);
thus the proof gives literal equality, not merely an equality modulo the
empty word. Both inclusions follow.
\end{proof}
\end{annexitem}

\begin{annexitem}{annex:p09.recombination-wrapper-separation}{thm:p09-recombination-wrapper-separation}{Extraction of the original language and parasitic wrappers}
\begin{proof}
The alphabets are disjoint at \(c,d\). Consequently
\(L_1\cap V^*=L'\cap V^*=L\), while
\(L_2\cap cV^*d=L'\cap cV^*d=cLd\).
Applying the declared restricted morphism to the latter gives \(L\).
The two closure hypotheses imply that membership of any of \(L_1,L_2,L'\)
in \(FA\) would imply \(L\in FA\), a contradiction.

The two wrapper types are disjoint. A word of \(L_1\) is either empty,
begins in \(V\), or begins with \(d\); its end is empty, in \(V\),
or \(c\). A word of \(L_2\) begins with \(c\) and ends with \(d\).
For every \(z\in L'\), both \(cz\notin L'\) and \(dz\notin L'\).
If \(z\in L_1\), then \(cz\) fails the required ending of \(L_2\)
and the beginning of \(L_1\). For \(dz\), the beginning is \(dv\)
with \(v\in V\), \(dd\), or just \(d\), none of which is a permitted
\(L_1\) beginning; its first letter also excludes \(L_2\).
If \(z\in L_2\), then \(cz\) begins \(cc\), excluding both types;
\(dz\) has final letter \(d\), excluding \(L_1\), but initial letter
\(d\), excluding \(L_2\). These checks include \(z=\lambda\).

Suppose \(L'=\sigma^*(L_0)\) with \(L_0\in FA\).
It is closed under every applicable rule of \(\sigma\).
Take any such application with first input
\(x=x_1u_1u_2x_2\in L'\), second input
\(y=y_1u_3u_4y_2\in L'\), and output \(z=x_1u_1u_4y_2\).
Closure puts \(z\in L'\).
If \(x\in L_1\), the word \(cxd\in L_2\) has the same occurrence
of the first recognition site. Splicing it with the same \(y\), at
that occurrence and its unchanged second site, yields \(cz\notin L'\).
If \(x\in L_2\), then \(dxc\in L_1\) gives \(dz\notin L'\).
Both contradict closure. This works even when a handle or site is empty,
and even when an output equals an original input.

Thus no rule is applicable on \(L'\) at all. Its iterated closure equals
its seed, so \(L'=L_0\in FA\), contradicting the first paragraph.
No hypothesis on the complexity of the rule set was used.
\end{proof}
\end{annexitem}

\begin{annexitem}{annex:p09.recombination-transition-coding}{thm:p09-recombination-transition-coding}{All transition paths from the finite no-triples seed}
\begin{proof}
A successful path consists of nonfinal transition letters whose target
state is the next letter's source, followed by one final letter, and begins
at \(S\). If there are \(m\) nonfinal letters, the no-triples condition
bounds a seed path's length by \(2m+1\). Over finite \(V\) there are
finitely many such words, and they can be enumerated and their path
conditions checked. The optional empty seed adds one word.
Every rule has one single-letter handle on each retained side and empty
other handles; hence its radius is at most one.

A rule \(A\#\$\#C\) cuts its first input immediately after \(A\)
and its second immediately before \(C\). The state's target in \(A\)
equals the source in \(C\). It therefore joins a successful path's
initial segment to a successful path's final segment with the correct
common state. No final-transition letter can be an \(A\), so a premature
final letter is not retained before the join. The output still begins at
\(S\) and ends in exactly one final letter. Induction proves that every
nonempty generated word is a successful path. The empty word participates
in no rule, because each operand needs its indicated transition letter.

Conversely use strong induction on successful path length. A path with no
tripled nonfinal letter is in the seed. Otherwise choose a nonfinal letter
\(A=[X,a,Y]\) occurring three times, and write
\(w=pAqArAs\). The words \(w'=pAqAs\) and \(w''=pArAs\) are successful
paths: the identical repeated letter has the same source and target,
so deleting the intervening portion joins matching states. Each is strictly
shorter than \(w\), even when \(q\) or \(r\) is empty, because a removed
copy of \(A\) has positive length. By induction both are available.

Let \(C\) be the first letter after the first \(A\) in \(w''\).
It exists, because another \(A\) follows, and its source is \(Y\).
The rule \(A\#\$\#C\) is present. Cut \(w'\) after its second
\(A\), retaining \(pAqA\), and cut \(w''\) before that \(C\),
retaining \(rAs\). Their splice is exactly \(pAqArAs=w\).
This also covers consecutive repetitions: then \(C=A\), possible only
when \(A\) is a loop, and that matching-state rule is present.
Thus all successful paths are generated.

Reading the middle component of each transition letter gives exactly the
terminal word of its grammar derivation, and every such derivation has its
successful transition path. Hence the coding gives the nonempty target
language. If it contains \(\lambda\), the optional seed supplies it;
otherwise no generated word codes to \(\lambda\). This proves literal
equality with the stated convention.
\end{proof}
\end{annexitem}

\begin{annexitem}{annex:p09.recombination-consuming-simulation}{thm:p09-recombination-consuming-simulation}{Single-copy phase control and complete consuming simulation}
\begin{proof}
Deleting identity grammar productions leaves the generated language
unchanged: remove those steps from every derivation. Thus \(u\ne v\).
All rule ranges are finite because the grammar itself is finite.
The reservoir support
and the initial support are finite. An infinite multiplicity is unchanged
by consuming or adding finitely many copies.

Write \(C(p,q)=X_1^2pYqX_2^2\), with \(p,q\in U^*\).
Its decoding is \(pq\). Initially the only active word is \(C(\lambda,S)\)
with multiplicity one. The two-letter left and right boundary blocks
ensure two available letters on each outer side of any selected occurrence.
Markers are distinct from \(U\).

\paragraph{A grammar step and its reassembly.}
If \(q=ut\), group 1 and reservoir \((r)v[r]\) give the two products
\[
 C(p,ut),\ (r)v[r]
 \quad\longmapsto\quad
 X_1^2pYu[r],\ (r)vtX_2^2.
\]
Choose \(\delta_1\delta_2\) as the two letters before \(Yu\) and
\(\beta_1\beta_2\) as the first two letters of \(tX_2^2\).
Their declared ranges are satisfied also for \(p=t=\lambda\).
Both products have multiplicity one; the previous active word is consumed.
Group 2 then gives
\[
 X_1^2pYu[r],\ (r)vtX_2^2
 \quad\longmapsto\quad
 C(p,vt),\ (r)u[r].
\]
Its \(\alpha\) is the first letter of \(tX_2^2\), hence in
\(U\cup\{X_2\}\). The decoded form changes from \(put\) to \(pvt\).
The two fragments are consumed; one active word remains.

\paragraph{Right cursor motion.}
For \(q=\alpha t\), group 3 with reservoir \(Z_1\alpha YZ_2\) yields
\[
 C(p,\alpha t),\ Z_1\alpha YZ_2
 \quad\longmapsto\quad
 X_1^2pY\alpha Z_2,\ Z_1\alpha YtX_2^2.
\]
Again the two exterior guards exist. Group 4, with \(\delta\) the last
letter of \(X_1^2p\) and \(\beta\) the first of \(tX_2^2\), yields
\[
 X_1^2pY\alpha Z_2,\ Z_1\alpha YtX_2^2
 \quad\longmapsto\quad
 C(p\alpha,t),\ Z_1Y\alpha Z_2.
\]
The second product is the opposite-moving reservoir, not another active
word. The decoding \(p\alpha t\) is unchanged.

\paragraph{Left cursor motion, including the rightmost gap.}
For \(p=s\alpha\), the repaired group 5 with \(Z_1Y\alpha Z_2\) gives
\[
 C(s\alpha,q),\ Z_1Y\alpha Z_2
 \quad\longmapsto\quad
 X_1^2s\alpha YZ_2,\ Z_1Y\alpha qX_2^2.
\]
For \(q\ne\lambda\), use the published three-letter guard: its first
letter is in \(U\), and the next two exist by the final \(X_2^2\).
For \(q=\lambda\), use the added exact \(X_2^2\) boundary rule.
This does not weaken the guard on other fragments. Group 6 then gives
\[
 X_1^2s\alpha YZ_2,\ Z_1Y\alpha qX_2^2
 \quad\longmapsto\quad
 C(s,\alpha q),\ Z_1\alpha YZ_2,
\]
using the last letter of \(X_1^2s\) and first of \(qX_2^2\).
The decoding is unchanged. Repeated left and right moves position \(Y\)
at any gap of a decoded word, including both endpoints.

Consequently every grammar derivation can be simulated: move \(Y\) just
before the chosen occurrence of its left side, perform the matching
cut/reassembly pair, and repeat. If the derived word is terminal, move
\(Y\) to the first gap. This proves the simulation direction once cleanup
is supplied below.

\paragraph{Excluding cross-combinations and false computations.}
The following exhaustive phase check proves the converse, not merely the
existence of the desired steps. Prior to cleanup, there is either one
completed \(C(p,q)\) or exactly the two multiplicity-one fragments of one
of the three displayed cut phases, in addition to reservoirs and grammar
waste words \((r)u[r]\). There cannot be fragments from two computations:
each cut consumes the sole completed active word, and each reassembly
consumes its two fragments.

Reservoirs cannot splice with one another. Groups 1 and 3 need two letters
from \(U\cup\{X_1\}\) before \(Y\), whereas the moving reservoirs have
only \(Z_1\) or \(Z_1\alpha\) there. Group 5 needs a predecessor
\(\delta\in U\cup\{X_1\}\) before \(\alpha Y\), excluding \(Z_1\).
Groups 2, 4 and 6 need their matching left fragments as first operands.
The \(YY\) reservoir has no \(X_1^2Y\) second site, and no initial
word has \(Y^3\), as needed for group 8.

In a grammar cut phase, the left fragment ends in \(Yu[r]\),
with \(1\leq|u|\leq2\).
It cannot supply the two-symbol right guards of groups 1 or 3,
because the bracketed tag \([r]\) is not in \(U\cup\{X_2\}\).
The original group-5 three-letter guard encounters that tag, while its
new boundary case requires \(X_2^2\), absent from this fragment.
The right fragment contains no \(Y\). Groups 4 and 6 require \(Z_2\)
and \(Z_1\), absent from these fragments. Only group 2 can rejoin them;
its tags force the same production label \(r\).
The reservoir \((r)v[r]\) cannot be a group-2 second operand, because
the required letter after \(v\) would be \([r]\), excluded from
\(\alpha\)'s range.

A waste \((r)u[r]\) cannot serve as a group-1 reservoir because
\(u\ne v\). It has no \(Y\), so is not a first operand.
It can serve as a group-2 second operand only if \(u\) begins with
\(v\alpha\) for an allowed \(\alpha\). With the matching left fragment
the two outputs are exactly that same left fragment and that same waste:
this is a stuttering event, not a new active word. No differently labeled
tag can match. Such events may repeat, but do not change the phase invariant.

In a right-moving cut phase the left fragment ends in \(Y\alpha Z_2\)
and the right starts in \(Z_1\alpha Y\).
Groups 1 and 3 require two allowed symbols after the site, which the
\(Z_2\) endpoint prevents on the left; the \(Z_1\) beginning prevents
their two allowed left guards on the right. Group 5 either lacks its
allowed predecessor or its two allowed right guard letters.
There are no grammar tags for group 2. Only group 4 matches the left
ending and right beginning, with the identical \(\alpha\).
The analogous check for the left-moving phase uses endings
\(\alpha YZ_2\) and beginning \(Z_1Y\alpha\), and permits only group 6.
A moving reservoir cannot replace a right fragment in a reassembly:
its next letter is \(Z_2\), not the allowed \(\beta\).

Group 7 requires \(X_1^2Y\) and a following window belonging to its
listed terminal/right-boundary set. It cannot match a split right fragment,
which has no \(X_1\), or a split left fragment, whose tag or \(Z_2\)
occurs before a complete allowed window. In a completed phase it requires
\(p=\lambda\). Groups 1--6 are therefore exactly the grammar/cursor
operations already analyzed until cleanup starts. Induction on actual
multiset events proves the invariant, including stuttering events.

\paragraph{Terminal cleanup and its complete boundary.}
For a terminal word \(t\), take the completed word \(X_1^2YtX_2^2\).
The initial segment of \(tX_2^2\) belongs to group 7's window set:
use \(X_2^2\) if \(|t|=0\), \(tX_2^2\) if \(|t|=1\) or 2,
and the first three letters of \(t\) otherwise.
Splicing with \(YY\) gives
\[
 YY,\ X_1^2YtX_2^2
 \quad\longmapsto\quad tX_2^2,\ X_1^2Y^3.
\]
Group 8 consumes these two words, using the \(Y^3\) occurrence in the
second, and gives \(t\) together with \(X_1^2Y^3X_2^2\).
Thus every generated terminal word, including \(\lambda\), is produced.

Group 7 can also start cleanup when only the initial window is terminal
and a later letter of the decoded form is a nonterminal. Its outputs still
have that entire decoded form, now without \(Y\), and a marked trash word.
No group 1--6 can process the unmarked decoded form: their \(Y\) or
tag sites are absent. Group 8 only deletes its final \(X_2^2\).
The trash \(X_1^2Y^3\), or \(X_1^2Y^3X_2^2\) after rejoining, has no
\(U\)-letter next to its \(Y\)'s, no \(Z_1,Z_2\), and no grammar tag.
It cannot match groups 1--6, nor the group-7 second window, which begins
with a further \(Y\). No fresh active word can be formed. There is only
one trash word carrying the sole \(X_1^2\) block, so it cannot supply
two occurrences for a self-splice. The remaining reservoirs and waste words
still cannot initiate a computation.

Consequently a terminal output can only be the fully terminal decoded
form of a completed active word. By the established invariant that form
is a grammar sentential form, hence lies in \(L(G)\). This proves the
opposite inclusion. Finally, if \(|u|\leq2\), \(|v|\leq2\), the longest
handle is \(\delta_1\delta_2Yu\), of length at most five; checking the
remaining listed handles gives at most four. The repaired group has no
longer handle. Thus the source-case radius bound survives the repair.
\end{proof}
\end{annexitem}

\begin{annexitem}{annex:p09.recombination-finite-population}{thm:p09-recombination-finite-population}{The actual finite state space and its effective edges}
\begin{proof}
Each event subtracts two input occurrences and adds two output occurrences,
including coincident inputs or outputs with their full multiplicities.
It therefore preserves \(\operatorname{cnt}(M)\).
For the four handles and factorizations, the two output lengths add to
\[
 |x_1|+|u_1|+|u_4|+|y_2|
 +|y_1|+|u_3|+|u_2|+|x_2|=|x|+|y|.
\]
Thus total weight is preserved too.

Every word in a multiset of weight \(W\) has length at most \(W\).
There are finitely many words of those lengths over finite \(V\), including
\(\lambda\). Each of their multiplicities is an integer between zero
and \(n\). Enumerate those finitely many vectors and retain exactly the
ones of count \(n\) and weight \(W\). This constructs the finite vertex
set; every reachable multiset belongs to it by conservation.

For finite \(R\), examine every pair of input occurrences at a vertex,
every rule and every matching factorization. There are finitely many.
Check the two-copy condition for equal inputs, compute both complementary
products, and update the full multiplicity vector. This constructs exactly
the outgoing edges. For decidable \(R\), each handle is a factor of a
word of length at most \(W\); enumerate the finitely many quadruples of
such handles and ask the given membership procedure whether the encoded
rule belongs to \(R\). Then use the same construction.

Finite graph reachability decides whether a proposed multiset is reachable.
An infinite evolution exists exactly when a directed cycle is reachable:
an infinite path repeats a vertex by finiteness, while a reached cycle can
be traversed repeatedly. Compute all reachable outdegrees and take their
maximum to decide the indicated nondeterminism measure. Parallel site/rule
descriptions yielding the same successor count as one edge for this
outdegree; if labeled-event counts are desired they are a different measure.
Without fixed count, the equal-weight vertex set is not finite, since
adding arbitrary finite multiplicities of \(\lambda\) changes no weight.
Without effective rule membership, the edge construction is not an
algorithm. Both necessary qualifications are therefore explicit.
\end{proof}
\end{annexitem}
